\documentclass{report}
\usepackage{amsmath,amssymb}
\setlength{\parindent}{0mm}
\setlength{\parskip}{1em}
\begin{document}
\begin{center}
\rule{6in}{1pt} \
{\large Ben O'Neill \\
{\bf Parallel in time multigrid for Nonlinear Problems}}

2210 Walnut St Apt 1 \\ Boulder CO 80302
\\
{\tt ben.oneill@colorado.edu}\\
Jacob Schroder\\
Thomas Manteuffel\end{center}

Standard sequential time marching schemes limit parallelism to the
spacial domain. With computer architectures growing in size rather than
clock speeds, additional speed-up must come from greater parallelism. We
propose a multigrid reduction method that incorporates temporal
parallelism into general time-stepping routines, allowing for dramatic
speed-ups on large architectures. For a nonlinear equation, each
iteration of the parallel-in-time method requires an expensive,
nonlinear, spatial solve. Using a simple nonlinear equation, with a
Picard method as the nonlinear solver, we investigate several methods for
reducing the computational cost of this spatial solve, including reducing
solver accuracy on coarser levels and introducing spatial coarsening.


\end{document}
